\chapter*{卷二\quad 岁计}
\phantomsection
\addcontentsline{toc}{chapter}{卷二\quad 岁计}
\markboth{卷二\quad 岁计}{卷二\quad 岁计}

\noindent
卷二十片，是全书中数字最密的一卷。它们记的不是某一年的账，是「一直如此」的那本账。

这一卷的释文最完整，校记却最难写。数字俱全，意思不明：每一片都能读出数，没有一片说清这个数是对什么而言的。

\begin{jian}{简十}
岁计口，凡三十二万一千。匠六万，工十一万，余为居者。

凡三口出一工。工者不入匠籍。匠者不入工籍。
\end{jian}
\chuchu{原件残片 37 号，长 20.5 cm，金属，字口较浅。}
\xiaojì{「三十二万一千」较简十四所称三十一万二千为多。此片字口较浅，与他片不类，疑为另一次清点，或为草稿；同卷之中两数并存而不相校，是他片所无。\par
「匠六万，工十一万」，两项相加十七万，余十四万一千为居者。三项之和三十一万一千，与本片所称三十二万一千差一万。此差不可解。\par
「凡三口出一工」，为征工之法。若以三口出一工计，则十一万工当出于三十三万人，与全数不合；可见「工」与「口」不是同一套编法。此点下编第六章另论。}

\begin{jian}{简十一}
匠籍别为一册。工籍别为一册。居者不入册，以口计。

册各有吏，不相知。
\end{jian}
\chuchu{原件残片 39 号，长 14.8 cm，陶瓷。}
\xiaojì{「居者不入册」，说明城中人口中最大的一部分不在任何名册上。此点对下编第六章的推算影响甚大：所有残存的数字，都是关于那十七万有籍者的，而十四万居者只有一个总数。\par
「册各有吏，不相知」，是一项极重要的制度。三册分掌，互不相通，则任何一名吏都不能掌握全城人口。此制与简八「板与册不同者以板为正」合观，可见城中对记录的控制是一以贯之的。}

\begin{jian}{简十二}
粟支九年，水支六年，气支三年。

三物之支，同时而计，不同时而尽。

尝有一岁，粟不足三月。减自上始，居者日一升，工如故。
\end{jian}
\chuchu{原件残片 41 号，长 9.4 cm，金属，为全书中最短一片。}
\xiaojì{三数并列，无主语。所支者为库存总量，抑或为每人每年之定额，不可辨。若为库存总量，则与简十四所记口数、简十五所记日给相合，可推出仓容；若为定额，则无解。本书取下编第六章所论之前一说。\par
「同时而计，不同时而尽」，为全卷最要紧的一句话。三物之中气最先尽，水次之，粟最后。若三者确为存量，则此城在理论上有一个明确的「断气之日」，而那个日子远在撤离之前。这一点与简二十七所记地热之衰相参，可见撤离不是等到粮尽，是等到气尽之前。\par
末段记一次粮荒：某一年，存粟只够三个月。处置是「减自上始」，居者每日一升，工匠照旧。\par
一升是定额的一半。按简十五「人日二升」计，居者的口粮减半，而工与匠不减。这与简十五「不足则减，减自上始」正合：减的顺序是先居者，再工，再匠。\par
这一条把那个「减」字落到了实处。它不是一条原则，是一次真的发生过的削减；被减的人占全城十四万一千居者中的大多数。\par
哪一年，本片未记。}

\begin{jian}{简十三}
新入者，先居外舍三月，乃入环。入环之日，给籍，给衣，给器。

器不继，以旧器补。旧器有记，记其旧主。
\end{jian}
\chuchu{原件残片 43 号，长 17.1 cm，金属。}
\xiaojì{「新入者」，指自地上归来之人。此类人在城中占比极小，然制度为其专设，可见往来未曾断绝。三月之期，与他处所记的隔离期皆合。\par
「旧器有记，记其旧主」，谓凡补给的旧器上保留原主之名。这一细节与全书的气氛相反：一个以「名可削」为制度的地方，却在器物上留着人的名字。此点下编第一章与第八章各有一说。}

\begin{figure}[htbp]
\centering
\begin{tikzpicture}[x=1mm,y=1mm]
  \useasboundingbox (-8,-16) rectangle (76,30);
  % 片身
  \fill[shade] (0,0) rectangle (56,18);
  \draw[ink,line width=0.5pt] (0,0) rectangle (56,18);
  % 行界与字界
  \foreach \x in {7,14,21,28,35,42,49}{\draw[rulegrey,line width=0.15pt] (\x,0) -- (\x,18);}
  \draw[rulegrey,line width=0.3pt] (0,9) -- (56,9);
  % 第一行：计口三十一万二千
  \foreach \i/\c in {1/计,2/口,3/三,4/十,5/一,6/万,7/二,8/千}{
    \node[font=\scriptsize] at (\i*7-3.5,13.5) {\c};}
  % 第二行：实居十一万□千
  \foreach \i/\c in {1/实,2/居,3/十,4/一,5/万,6/九}{
    \node[font=\scriptsize] at (\i*7-3.5,4.5) {\c};}
  \node[font=\scriptsize] at (49-3.5,4.5) {千};
  % 蚀去（不在其上压线，只填浅灰并以细边标出）
  \fill[faintgrey] (42,0) rectangle (49,9);
  \draw[rulegrey,line width=0.15pt] (42,0) rectangle (49,9);
  \node[font=\scriptsize,text=pencil] at (49-3.5,4.5) {□};
  % 尺寸
  \draw[<->,rulegrey,line width=0.3pt] (0,-3) -- (56,-3);
  \node[font=\scriptsize,text=pencil,anchor=north] at (28,-4) {18.4 cm（原件实测）};
  \draw[<->,rulegrey,line width=0.3pt] (60,0) -- (60,18);
  \node[font=\scriptsize,text=pencil,anchor=west,align=left] at (61,9) {1.7 cm\\两行};
  \node[font=\scriptsize,text=pencil,anchor=south west] at (0,18.6) {残片 26 号};
\end{tikzpicture}
\caption{简十四（残片 26 号）形制与刻写行款复原}
\label{fig:slip}
\figsource{据原件照片与一九六三年编目卡重绘。灰色一格为蚀去之处，即校记所补「百」字的位置。}
\end{figure}

\begin{jian}{简十四}
计口三十一万二千，实居十一万九千。

粟支九年，水支六年，气支三年。
\end{jian}
\chuchu{原件残片 26 号（编号与卷一次序不合，系一九六三年二次编目所致），长 18.4 cm，金属，正面有朱色残痕。}
\xiaojì{「十一万九千」，胶片本作「十一万九千六百」，残片此处蚀去，据文例补「百」字，仍存疑。\par
此片为全卷最重要一片：额定与实居并列，差额十二万一千，即空置近四成。第二行三数与简十二全同，可知简十二即此片之下半，二者或为同一条刻于两片者。\par
若此说成立，则简十二的「同时而计，不同时而尽」一句，当接在此片之后。}

\begin{jian}{简十五}
月给粟，人日二升。匠倍之。外工又倍之。

不足则减，减自上始。
\end{jian}
\chuchu{原件残片 45 号，长 13.6 cm，金属。}
\xiaojì{按二升计，十一万九千人日需二千三百八十石，岁需八十六万八千石。与简十二「粟支九年」相权，则仓容当在七百万石以上。此数为下编第六章推算之基础。\par
「不足则减，减自上始」，为分配原则。减自上始者，先减居者，后减工，后减匠。此句极短，却是全卷唯一一处涉及分配次序的记载。}

\begin{jian}{简十六}
工籍：捣者三千，负者一千二百，修者八百，余为杂役。

捣者不入环，月给别。
\end{jian}
\chuchu{原件残片 47 号，长 16.9 cm，陶瓷。}
\xiaojì{「捣者三千」，三千之数为全书所仅见，与他处工匠之数皆不合。「捣」字在此处无宾语，所捣为何物，本片未言。\par
「捣者不入环，月给别」，谓此三千人不入居住环，另给口粮。此种安排与简五「外工」之制相类而不同：外工居垣外，捣者不入环而未言居处。二者是否为同一批人，无据可断。}

\begin{jian}{简十七}
定额：人日三合，多者赏，少者罚。罚不过三，过三则移籍。

赏以粟，罚以役。

十四年，捣者十一人连日不足，罚役三日，不移籍。其一人再犯，移籍。
\end{jian}
\chuchu{原件残片 49 号，长 12.3 cm，金属。}
\xiaojì{「三合」为容积单位，其量当指成品。此片与简十六相接，可断为同一册。\par
「赏以粟，罚以役」，谓奖惩一律以劳役与口粮结算，不涉他物。此种安排在残片中仅此一见，可见城中的激励机制极为单一。\par
「过三则移籍」，移籍即改属他册。简十一言「册各有吏，不相知」，则移籍是一次跨系统的调动，非一名吏所能独断。移往何籍，本片未言。\par
末段是一次具体的处分：十一个人连着几天没做够定额，罚做三天工，不改名籍；其中一个又犯，改了籍。\par
这是全卷唯一一处能看出「罚不过三」如何执行的记载。十一人同时不足，说明那一天不是个人的问题，是定额本身出了问题；但处分照旧，只有「其一人再犯」被移籍。\par
为什么是这一人，本片未言。}

\begin{jian}{简十八}
凡籍，岁一易。易籍之日，旧籍焚。焚者二人，见者三人，皆署名。

焚余之灰，埋于垣外。

十九年，一册不焚。焚者自陈，见者三人同辞。罚焚者，册终不焚。
\end{jian}
\chuchu{原件残片 51 号，长 15.8 cm，金属。}
\xiaojì{旧籍岁焚，是本书最大的一个空洞：一九六三年所得者，皆为焚余或未及焚者，非全帙。凡下编所论，皆当以此为前提。\par
「焚者二人，见者三人，皆署名」，为焚籍之程序。署名者五人，则每一次焚籍都留下五个名字；而本书三十七片中，一个这样的名字也没有。可见所焚之籍与被焚之签，皆不在此。\par
「焚余之灰，埋于垣外」，与简六「限外百步不得置物」相冲。二禁相权，当是葬灰为例外。\par
末段记一桩事故：某一年，有一册没有烧。烧的人自己报了，在场的三个人说法一致。结果是烧的人受罚，那一册最终也没有烧掉。\par
「终不焚」三字没有解释。可以理解为没来得及补烧，也可以理解为有意留着。本片不偏任何一边。\par
这一条与本书的成立直接相关。若所有旧籍都照章焚毁，一九六三年能得到的只会是灰。我们今天能读到三十七片，是因为在这个制度里，总有一两次「不焚」。}

\begin{jian}{简十九}
事故之数：兔乱一万一千，他故八千，逃亡三千。凡二万二千。

逃亡者，籍除，名削。
\end{jian}
\chuchu{原件残片 53 号，长 19.2 cm，金属，右端有灼痕。}
\xiaojì{「兔乱一万一千」，「兔」字全卷仅两见，此为其一。「他故」「逃亡」二项，此前未见著录，为本书新增材料。\par
三项合计二万二千，与三十一万二千至十二万之差额三万中余八千之数，尚不能尽合，详下编第六章。\par
「逃亡者，籍除，名削」，为处置办法。逃者既除籍削名，则其不在任何一册之内；此点与简二十五「无名者不得登」相接，是全书第二处大矛盾。}
